\section{IX: required index demands}
\label{sec:ix}
Between RA and physical selection, IX records which attribute-keyed accesses
the fixed plan demands. Physical derives uses from its operands. IX performs
no data-structure selection and emits no runtime construction. The compiled
chain continues
\[
 \mathrm{RA}\longrightarrow\mathrm{IX}
 \xrightarrow{\text{lower}}\mathrm{Physical}
 \xrightarrow{\text{select}}\mathrm{Structures}
 \longrightarrow\mathrm{Rust},
\]
The caller selects lowering, optional rewrites, selection, and printing as
public macro paths in the checked pipeline. Their arguments are local to each
pass; the complete target IR records every resulting decision.

\paragraph{Boundary syntax.}
Let $r,l$ name exact RA versions, $j$ an operator occurrence, $i$ an index,
and $C$ an ordered list of logical column ordinals. For a binary stream demand, $r$ is the right operand of a join, semijoin
or anti; a const/range demand names the exact select input.
\begin{lstlisting}
IX  ::= RAProgram indexes { Def* }
Def ::= i = @r[C] (stream | const | range);
\end{lstlisting}
IX retains RA unchanged. The listing is the fragment used below; the complete
grammar is in \path{languages/cq-lang-ix/src/lib.rs}.

\paragraph{IX meaning and construction.}
Write IX as $(A,D)$, with $D:i\mapsto(r,C,supply)$.
Well-formedness requires well-formed RA, distinct index identities and
descriptors, and strictly increasing, in-range logical ordinals over exact RA
results. Every demand matches an RA join/semijoin stream read, an anti stream read
on its explicit free key, a constant/parameter point
conjunction, or constant order bounds on one column of an exact select input. Partial demand sets are locally valid; the analysis
correspondence additionally requires the complete sequence in first-use order
and the names \code{i0}, \code{i1}, etc., as independently checked in
\path{passes/cq-index-policy/src/correspondence.rs}.
Erasure preserves meaning: $\sem{A,D}_{IX}=\sem A_{RA}$.
Traverse the fixed RA definitions in order, starting with empty $D$:
\begin{mathpar}
\Rule{IX-Reuse}{D(i)=(r,C,supply)}{D\vdash(r,C,supply)\Downarrow(i,D)}

\Rule{IX-New}{(r,C,supply)\notin\operatorname{ran}D\quad i=i_{|D|}}
 {D\vdash(r,C,supply)\Downarrow(i,D[i\mapsto(r,C,supply)])}

\Rule{IX-Join}{o\in\{\bowtie,\ltimes\}\quad
 C=[c\mid 0\leq c<|H(r)|,\ H(r)_c\in H(l)]_{\rm increasing}\\
 D\vdash(r,C,stream)\Downarrow(i,D')}
 {D\vdash j=l\mathbin{o}r\Downarrow D'}

\Rule{IX-Const}{\theta\text{ is a point conjunction on columns }K\\
 \text{each term is a constant equality or literal membership}\\
 C=\operatorname{ordinals}_{H(r)}(K)\quad
 D\vdash(r,C,const)\Downarrow(i,D')}
 {D\vdash j=\sigma_\theta r\Downarrow D'}

\Rule{IX-Range}{\theta\text{ conjunction of constant/parameter order bounds on }c\\
 C=[\operatorname{ordinal}_{H(r)}(c)]\quad
 D\vdash(r,C,range)\Downarrow(i,D')}
 {D\vdash j=\sigma_\theta r\Downarrow D'}

\Rule{IX-Other}{d\in\{\mathsf{Unit,Rename,Project,Guard,Bind,Union}\}\\
 \text{ or a select matching neither IX-Const nor IX-Range}}
 {D\vdash d\Downarrow D}
\end{mathpar}
For anti, IX-Join uses only the written key $F$, not the wider heading
intersection: $C=\operatorname{ordinals}_{H(r)}(F)$.
Fresh $i_n$ prints as \code{i0}, \code{i1}, etc. Together with the anti case, these cover the current RA operator enum.
$C=[]$ is the unit tuple for the admitted Cartesian case. Physical lowering
derives operands from RA and checks probe/semi/anti, point and range supplies in $D$; merge has separate
ordered-read obligations and its left side is excluded from IX demand checking;
this demand-membership check belongs to the IX-to-Physical pass. The standalone
Physical contract derives local headings, reads and modes without IX. Fine
plans elaborate to RA operators covered by these same demand rules; no WCOJ
complexity guarantee follows.

\paragraph{Physical and Structures.}
\begin{lstlisting}
Physical ::= physical { Decl* Def* return [canon] r as Head; }
Def ::= r = (pipe | let | exists) [canon] Op;
Op ::= Unit | Rename | select ... | point select ... | range select ...
     | Bind | Guard | Project | merge [K] l, r
     | probe [first|smallest] [K] l, r
     | semi [eager|lazy] [K] l, r | anti [eager|lazy] [F] l, r
     | union [any|all] [fused] r1, r2, ... [BlockOrigin]
Structures ::= Physical structures { Entry* }
Entry ::= r = Type [K]... [payload[P] constants{c = term, ...}];
        | r host reader = Type [K]...;
        | r rows = Type; | r canon = Type; | r head = Type;
\end{lstlisting}

Physical restates every IX result with one physical operator (probe, merge,
semi, range, or the logical operator's own name) and the mode its reads allow
(pipe, let, exists), marking canonical results. Structures records concrete
built types and key levels, per-read host fallbacks, retained row collections,
and canonical-boundary and computed-head dedup types. Its local contract
discharges the recorded plan's read and uniqueness premises.
Host renames use type dispatch without capability declarations. Below Structures a printer emits one pull program
over the \code{cq-storage} operation traits. Their local contract rules and
pass correspondence judgments appear in Section~\ref{sec:pass-contracts}.
Complete grammars and representation obligations also appear in \path{languages/cq-lang-physical/src/lib.rs},
\path{languages/cq-lang-structures/src/lib.rs}. The lowering pass checks RA
erasure and exact IX demands; default selection copies Physical, while its
correspondence check permits the Physical rewrite relation and checks the
target contract. Neither independent check reruns its constructor.

\subsection{Recorded block execution and representations}
Union forwards its operands' scans. Anti requests a membership demand on the
exact free heading $F$, independently of wider common columns. Physical records
\code{union [any|all] [fused]} (witness choice only for a test interface),
with a definition-level \code{canon} or deferred canonicalization, and
\code{anti eager|lazy [F] r, s}. A plain union defers dedup to a final canon
boundary; a proven-disjoint union omits cross-arm canon. Structures records
the dedup provider and can build a derived-relation index directly from the
chained scan using its filtered-index representation.
\begin{mathpar}
\Rule{PH-Block}{\text{branch scopes derived from union origins and fixed operands}\\
 \text{ancestor/local source-before definitions remain before each Rust node}}
 {\mathsf{blockBeforeOK}(source,target)}
\end{mathpar}
Defaults use written branch order, test union \code{any}, binding union
\code{canon} unless differing interface literals prove disjointness, and a
fused prefix stream for per-row arms. In a fused group all arm operands can
read the same prefix pipe; execution still streams that prefix once. Core pull realization and the external push printer consume these recorded
choices. \code{any} records a witness stop (per prefix row for a fused test,
or for a unit heading); hoisted free-heading results must retain every free
tuple. \code{all} omits this \code{take(1)} bound, allowing all arm rows to
stream. It denotes union nonemptiness, not conjunctive truth; downstream
demand may still stop consumption. Fallback uses
if-empty control flow and never evaluates a later Rust arm after a witness.
